\subsection{The criterion as bounds on an infimum}

\cite[Chapter~3, Proposition~3.2]{Wolf2012Quantum} states the spectral
criterion for $n$-positivity as two inequalities for
$\inf_\psi\bra{\psi}\tau\ket{\psi}$, the infimum being taken over the
normalized vectors of Schmidt rank $n$. Those vectors are the ones obtained by
embedding $\C^n$ in the second factor $\C^{D}$, so they are the vectors of
Schmidt rank at most $n$: this is the set on which $n$-positivity of a
Hermitian map is tested, and the set whose expectations are collected below.

\begin{definition}[Expectations at bounded Schmidt rank]
    \label{def:schmidt_rank_le_expectations}
    \lean{Matrix.schmidtRankLEExpectations}
    \leanok
    \uses{def:schmidt_rank}
    For a matrix $\tau$ on $\C^{D'}\otimes\C^{D}$ and a bound $n$, the set
    \begin{align}
        E_n(\tau)
         & =\{\Re\bra{\psi}\tau\ket{\psi}:\|\psi\|=1,\ \SR(\psi)\le n\}
        \subseteq\R
        \label{eq:channel_schmidt_rank_le_expectations}
    \end{align}
    collects the real parts of the quadratic forms of $\tau$ in the normalized
    vectors of Schmidt rank at most $n$. For Hermitian $\tau$, the setting of
    every statement below, the quadratic form is real and $E_n(\tau)$ is the
    set of expectations of $\tau$ in those vectors.
\end{definition}

\begin{lemma}[The expectation set has an infimum]
    \label{lem:schmidt_rank_le_expectations_inf_exists}
    \lean{Matrix.schmidtRankLEExpectations_nonempty,
        Matrix.IsHermitian.schmidtRankLEExpectations_bddBelow,
        Matrix.mem_schmidtRankLEExpectations}
    \leanok
    \uses{def:schmidt_rank_le_expectations}
    Assume $D',D\ge1$. For $n\ge1$ the set $E_n(\tau)$ is nonempty, and for
    Hermitian $\tau$ it is bounded below by the least eigenvalue $\nu_{\min}$ of
    $\tau$; the expectation of $\tau$ in any normalized vector of Schmidt rank at
    most $n$ is one of its elements. Hence $\inf E_n(\tau)$ exists.
\end{lemma}

\begin{proof}
    \leanok
    \uses{thm:schmidt_rank_basic_bounds, thm:spectral_bound_schmidt_rank_top}
    A normalized product vector has Schmidt rank one, hence Schmidt rank at
    most $n$, so $E_n(\tau)$ is nonempty. For the lower bound,
    Theorem~\ref{thm:spectral_bound_schmidt_rank_top} gives
    $\nu_{\min}\le\bra{\psi}\tau\ket{\psi}$ for every normalized $\psi$,
    without any Schmidt-rank restriction.
\end{proof}

\begin{theorem}[Lower bound on the Schmidt-rank infimum]
    \label{thm:spectral_infimum_lower_bound}
    \lean{Matrix.IsHermitian.le_sInf_schmidtRankLEExpectations}
    \leanok
    \uses{def:schmidt_rank_le_expectations, def:ky_fan_norm,
        def:reduced_eig_density}
    Let $\tau$ be a Hermitian operator on $\C^{D'}\otimes\C^{D}$ with
    eigenvalues $\nu_i$, normalized eigenvectors $\phi_i$ and reduced densities
    $\rho_i=\tr_2\ketbra{\phi_i}{\phi_i}$. Assume $D',D\ge1$, let $1\le n<D'$
    and let $\nu_0\ge0$ be a lower bound for the positive eigenvalues, the smallest
    positive eigenvalue in \cite[Chapter~3]{Wolf2012Quantum}. Then
    \begin{align}
        \nu_0+\sum_{i:\nu_i\le0}(\nu_i-\nu_0)\|\rho_i\|_{(n)}
         & \le\inf E_n(\tau).
        \label{eq:channel_spectral_infimum_lower}
    \end{align}
    This is \cite[Chapter~3, equation~(3.7)]{Wolf2012Quantum} for $n<D'$. The
    indices $n\ge D'$ are Theorem~\ref{thm:spectral_infimum_top}, and the two
    together give the equation at every $n\ge1$.
\end{theorem}

\begin{proof}
    \leanok
    \uses{thm:spectral_lower_bound_schmidt_rank,
        lem:schmidt_rank_le_expectations_inf_exists}
    Theorem~\ref{thm:spectral_lower_bound_schmidt_rank} bounds
    $\bra{\psi}\tau\ket{\psi}$ below by the left-hand side at every normalized
    $\psi$ of Schmidt rank at most $n$, that is, at every element of
    $E_n(\tau)$. A lower bound for every element of a nonempty set is a lower
    bound for its infimum.
\end{proof}

\begin{theorem}[Upper bound on the Schmidt-rank infimum]
    \label{thm:spectral_infimum_upper_bound}
    \lean{Matrix.IsHermitian.sInf_schmidtRankLEExpectations_le}
    \leanok
    \uses{def:schmidt_rank_le_expectations, def:ky_fan_norm,
        def:reduced_eig_density}
    With $\tau$, $\phi_i$ and $\rho_i$ as above and $D',D\ge1$, let
    $1\le n<D'$ and suppose every eigenvalue of $\tau$ is at most $\nu$. Then,
    for every eigenvector index $j$,
    \begin{align}
        \inf E_n(\tau)
         & \le\nu+(\nu_j-\nu)\|\rho_j\|_{(n)}.
        \label{eq:channel_spectral_infimum_upper}
    \end{align}
    Taking $\nu$ to be the largest positive eigenvalue and $j$ the index of the
    unique non-positive eigenvalue $\nu_-$, which is the situation in which all
    other eigenvalues are strictly positive, gives
    $\inf E_n(\tau)\le\nu+(\nu_--\nu)\|\rho_-\|_{(n)}$, that is,
    \cite[Chapter~3, equation~(3.8)]{Wolf2012Quantum} for $n<D'$. The indices
    $n\ge D'$ are Theorem~\ref{thm:spectral_infimum_top}, and the two together
    give the equation at every $n\ge1$.
\end{theorem}

\begin{proof}
    \leanok
    \uses{thm:spectral_upper_bound_schmidt_rank,
        lem:schmidt_rank_le_expectations_inf_exists}
    Theorem~\ref{thm:spectral_upper_bound_schmidt_rank} supplies a normalized
    $\psi$ of Schmidt rank at most $n$ whose expectation is at most the
    right-hand side. Its expectation lies in $E_n(\tau)$, so the infimum of
    $E_n(\tau)$ is at most that expectation, hence at most the right-hand side.
\end{proof}

\begin{lemma}[The Ky-Fan norm of a reduced eigenvector density at full index]
    \label{lem:reduced_eig_density_ky_fan_top}
    \lean{Matrix.IsHermitian.trace_reducedEigDensity,
        Matrix.IsHermitian.kyFanNorm_reducedEigDensity_eq_one}
    \leanok
    \uses{def:ky_fan_norm, def:reduced_eig_density}
    Each reduced eigenvector density has unit trace, and for $n\ge D'$ its
    Ky-Fan $n$-norm is that trace:
    \begin{align}
        \tr\rho_i &= 1,
        &
        \|\rho_i\|_{(n)} &= 1
        .
        \notag
    \end{align}
\end{lemma}

\begin{proof}
    \leanok
    The partial trace preserves the trace, so
    $\tr\rho_i=\tr\ketbra{\phi_i}{\phi_i}=\braket{\phi_i}{\phi_i}=1$. The Ky-Fan
    norm sums the $n$ largest eigenvalues in descending order, and an index past
    the dimension $D'$ contributes zero, so the sum stops growing at $n=D'$,
    where it is the sum of all eigenvalues of $\rho_i$, that is $\tr\rho_i$.
\end{proof}

\begin{theorem}[The Schmidt-rank infimum at the top index]
    \label{thm:spectral_infimum_top}
    \lean{Matrix.IsHermitian.sInf_schmidtRankLEExpectations_top,
        Matrix.IsHermitian.le_sInf_schmidtRankLEExpectations_top}
    \leanok
    \uses{def:schmidt_rank_le_expectations, def:ky_fan_norm,
        def:reduced_eig_density}
    Assume $D',D\ge1$. Once the Schmidt-rank bound reaches the dimension $D'$ of
    the first tensor factor the constraint is vacuous, and each Ky-Fan norm
    collapses to
    $\|\rho_i\|_{(n)}=\tr\rho_i=1$. For $n\ge D'$,
    \begin{align}
        \inf E_n(\tau)
         & =\nu_{\min},
        \label{eq:channel_spectral_infimum_top}
    \end{align}
    which is the right-hand side of~(3.8) there, since
    $\nu+(\nu_--\nu)\|\rho_-\|_{(n)}=\nu_-=\nu_{\min}$ when $\nu_-$ is the only
    non-positive eigenvalue. The right-hand side of~(3.7) reads
    $\nu_0+\sum_{i:\nu_i\le0}(\nu_i-\nu_0)$, which is a strictly weaker bound as
    soon as $\tau$ has two non-positive eigenvalues: for the spectrum
    $(-2,-1,3)$ it is $-6$ while $\nu_{\min}=-2$. It is still a lower bound: for
    those same $n$, and for $\nu_0\ge0$ a lower bound for the positive
    eigenvalues as in Theorem~\ref{thm:spectral_infimum_lower_bound},
    \begin{align}
        \nu_0+\sum_{i:\nu_i\le0}(\nu_i-\nu_0)\|\rho_i\|_{(n)}
         & \le\inf E_n(\tau),
        \label{eq:channel_spectral_infimum_top_lower}
    \end{align}
    and~(3.7) holds at those indices as well. With
    Theorems~\ref{thm:spectral_infimum_lower_bound}
    and~\ref{thm:spectral_infimum_upper_bound} this covers every $n\ge1$, hence
    the range $1\le n\le D$ of
    \cite[Chapter~3, Proposition~3.1]{Wolf2012Quantum}, where $D$ is the
    dimension of the second tensor factor.
\end{theorem}

\begin{proof}
    \leanok
    \uses{thm:spectral_bound_schmidt_rank_top,
        thm:schmidt_rank_basic_bounds,
        lem:schmidt_rank_le_expectations_inf_exists,
        lem:reduced_eig_density_ky_fan_top}
    Every vector has Schmidt rank at most $D'$, so for $n\ge D'$ the set
    $E_n(\tau)$ is the set of expectations in all normalized vectors. By
    Theorem~\ref{thm:spectral_bound_schmidt_rank_top} that set is bounded below
    by $\nu_{\min}$ and contains $\nu_{\min}$, attained at an eigenvector for
    the least eigenvalue. This gives~(\ref{eq:channel_spectral_infimum_top}).
    For~(\ref{eq:channel_spectral_infimum_top_lower}), let $j$ be a minimizing
    index. If $\nu_{\min}\le0$ then $j$ belongs to the summation range and
    \begin{align}
        \nu_0+\sum_{i:\nu_i\le0}(\nu_i-\nu_0)
        &= \nu_{\min}+\sum_{i:\nu_i\le0,\ i\ne j}(\nu_i-\nu_0)
        \le \nu_{\min},
        \notag
    \end{align}
    every remaining summand being nonpositive. If instead every eigenvalue is
    positive the summation range is empty and the left-hand side is
    $\nu_0\le\nu_{\min}$.
\end{proof}
